\printedpage{44}
% head: A. Grothendieck : Dérivateurs, Ch. XV, Édition des Universités de Montpellier II et Paris VII

\msp{page 80}
small set, because of the hypotheses made on $\M$), the ordinals $\alpha$ for which $A_\alpha$ exists satisfy $\mathrm{card}\,A_\alpha \leq \pi$. So there must exist an $\alpha$ such that $A_\alpha = B$. But then $i_\alpha : A \longrightarrow A_\alpha$ is equal to $f$, hence $f \in \widetilde{C^{\pi}}$, \quad \textit{q.e.d.}

It remains then to prove 6.4. For this, let $I$ be the set of subobjects of $B$ which are in $\Filt^{\pi}$, ordered by inclusion. (If $B$ is in $\Filt^{\pi}$, $A$ too, so $f$ is already in $C^{\pi} \subset \widetilde{C^{\pi}}$, and there is nothing to prove.) It is a filtering set large relative to $\pi$ (and even $\pi$-\textit{adapted}), whose $\varinjlim$ is $B$ (5.3, p. 46). Consider the
\[
A_i \;=\; B_i \cap A \;,
\]
they are in $\Filt^{\pi}$, as subobjects of objects of $\Filt^{\pi}$ (condition $(d)$ on the cardinal filtration). Thus, the arrow $f$ appears as a filtering direct limit of the arrows
\[
f_i \;:\; A_i \longrightarrow B_i \;.
\]
It is trivial that the system of the $A_i$ is cofinal in $\Ss_{\pi}(A)$, hence $A = \varinjlim A_i$, whence $f = \varinjlim f_i$.

\msp{page 81}
On the other hand, I say that the direct system of the $f_i$ in $\underline{\Fl}(\M)$ is \textit{well adapted to $\pi$.} We already know that $I$ is well adapted to $\pi$, the sup in $I$ of a filtering subset $J$ of cardinal $\leq \pi$ being the sup in the ordered set $\Ss(B)$ (which sup is in $\Ss_{\pi}(B)$). It is essential, for our argument, to know that if $i_J = \sup J$, then not only $B_{i_J} = \varinjlim\limits_{i\in J} B_i$, but \marginpar{$\|$}% marginal double vertical bar printed in the left margin at this line
also $A_{i_J} = \varinjlim\limits_{i\in J} A_i$, so that the direct system of the $f_i$ is $\pi$-adapted in $\underline{\Fl}(\M)$. It seems that we therefore need an exactness condition on $\M$ close to $(e)$, namely:

\begin{itemize}
\item[\cercle{e$'$}] The filtering $\varinjlim$ functors commute with base change.
\end{itemize}

(It suffices it here in the very particular case of a base change by a monomorphism $A \longrightarrow B$, and of a filtering direct limit of subobjects of $B$ with inclusion morphisms as transitions.) Let us note that \cercle{e}, \cercle{e$'$} generally arise as a consequence of the condition % [sic] Il le suffit -> Il suffit, ungrammatical as printed

\begin{itemize}
\item[(e$_0$)] The filtering $\varinjlim$ in $I$ are left exact.
\end{itemize}

\msp{page 82}
With this laid down, we thus see, $C$ being $\pi$-accessible and $I$ large relative to $\pi$, that the set $I_C$ of $i \in I$ such that $f_i \in C$ is \textit{cofinal} in $I$. But since $f_i \in \Filt^{\pi}$, $f_i \in C$ is equivalent to $f_i \in C^{\pi}$. Since $A \longrightarrow B$ is not an isomorphism, there exists an $i_0$ such that $B_{i_0} \not\subset A$, then an $i \geq i_0$ such that $f_i \in C^{\pi}$, and also $B_i \not\subset A$ (since $B_i \supset B_{i_0}$). Thus we have proved this:

\begin{itemize}
\item[($*$)] Under the hypotheses of (6.4) and with the hypothesis \cercle{e$'$} on $\M$, there exists a subobject $B'$ of $B$ satisfying the following conditions.
\begin{itemize}
\item[\textit{a})] $B' \in \Filt^{\pi}$.
\item[\textit{b})] $B' \not\subset A$.
\item[\textit{c})] Setting $A' = B' \cap A$, the inclusion
\[
i' \;:\; A' \;\longrightarrow\; B'
\]
is in $C^{\pi}$.
\end{itemize}

\printedpage{45}
% head: A. Grothendieck : Dérivateurs, Ch. XV, Édition des Universités de Montpellier II et Paris VII

Moreover, the set of $J \subset \Ss_{\pi}(B)$ formed by the $B'$ in question is cofinal in $\Ss_{\pi}(B)$, and even $\pi$-cofinal.
\end{itemize}

We still have to prove this last point -- that $J$ (which we know is cofinal) is even $\pi$-cofinal, \textit{i.e.} stable under sup (in $\Ss_{\pi}(B)$) of filtering subsets $J' \subset J$ such that $\mathrm{card}\, J' \leq \pi$. But this follows at once from the fact that $C$ being stable under

\msp{page 83}
filtering direct limits, $C^{\pi} = C \cap \Filt^{\pi}$ is stable under filtering $\varinjlim$ of cardinal $\leq \pi$, by virtue of condition $(d)$ on $\Filt^{\pi}$. Consider now the cocartesian square
\[
\begin{tikzcd}
A' \arrow[r, hook, "i'"] \arrow[d, hook] & B' \arrow[d] & & \\
A \arrow[r, "i"] & C' = A \amalg_{A'} B' \arrow[r, "j"] & B \;.
\end{tikzcd}
\]
Since $i' \in C^{\pi}$, we have $i \in \widetilde{C^{\pi}}$ (for $\widetilde{C^{\pi}}$ stable under cobase change). On the other hand, I need the fact that $j$ is in $C$, and I do not even know that it is a monomorphism! So we have to go back to the indexed notation
\[
\begin{tikzcd}
A_i \arrow[r, hook, "\alpha_i"] \arrow[d, hook] & B_i \arrow[d] & \\
A \arrow[r, hook, "\beta_i"] & C_i \arrow[r, "\gamma_i"] & B \;.
\end{tikzcd}
\]
Since the formation of amalgamated sums commutes with filtering direct limits, we see that $\varinjlim C_i \simeq A \amalg_{\varinjlim A_i} \varinjlim B_i = A \amalg_A B = B$, hence the direct limit of the $\gamma_i$ is an isomorphism (essentially $\mathrm{id}_B$). On the other hand, the same argument shows that the direct system of the $C_i$ $(i \in J)$ in $\M$ is $\pi$-adapted, hence also that of the $\gamma_i$\footnote[28]{\textit{Except} that we do not know that the $C_i$ are in $\Filt^{\pi}$ -- in fact, they will not be in general, since we do not assume that $A, B \in \Filt^{\pi}$.}
\msp{page 84}
in $\underline{\Fl}(\M)$. Since $\varinjlim \gamma_i$ is an isomorphism, it is in $C$. Since $C$ is $\pi$-accessible, it follows that the set $J'$ formed by the $i \in J$ such that $\gamma_i \in C$ is cofinal as well\footnote[29]{\textit{Botched, cf.} preceding annotation.}. So we find $i$ with $\alpha_i \in C^{\pi}$, $\gamma_i \in C$, which establishes lemma 6.4, and completes the proof of theorem 6.1.

\section*{7\quad Construction of closed Quillen triples with $\boldsymbol{C} = \mathbf{Mon} \subset \mathbf{Cof}_{\boldsymbol{W}}$}

Finally, all the beautiful developments of paragraphs 5, 6 do not give me the desired ``sweeping result''. Perhaps it is false -- perhaps the property of generation
\printedpage{46}
% head: A. Grothendieck : Dérivateurs, Ch. XV, Édition des Universités de Montpellier II et Paris VII
``$^{\sim}$'' by a small subset (for accessible, or L-accessible sets) is satisfied much more rarely, than I would have imagined. To do things properly, I would have to work through to the % resolved: the quoted raised wave denotes Q-closure; see corrections.csv, leaf 13.
\msp{page 85}
end a few examples, e.g. that of $\Cof_W$, and of $W^{\mathrm{univ}}$ in Cat, for a more or less general basic localizer, and especially in the case $W = \mathbf{W}_{\infty}$. Surely $\Cof_W$ and $W^{\mathrm{univ}}$ are accessible (I ought to check it nevertheless), and if the hypothetical ``theorem'' of paragraph 6 is false for them, at least I will know for sure. But there still remains the question: if $C = \widetilde{C}$ is stable under filtering $\varinjlim$, and of the form $\widetilde{C_0}$, $C_0$ small, and if $W$ is L-accessible as well, then is $C \cap W$ also of the form $\widetilde{TC_0}$, with $TC_0$ small? I have just tried once again to prove it, assuming moreover $C \subset$ mono, and see no way of getting there. So a counterexample would be needed to dispel these perplexities! And the best would be in Cat, while we're at it.

There nevertheless remains afloat a tangible positive result, which I will now try to pin down properly and prove:

\msp{page 86}
\textbf{Theorem 7.1.} \textit{Let $\M$ be a category stable under small filtering $\varinjlim$, $M_0 \subset \Ob\,\M$ a subset of} $\Ob\,\M$ \textit{generating by strict monomorphisms, more precisely we assume this:}

\begin{quote}
\itshape
For every monomorphism $i : A \hookrightarrow B$ in $\M$ which is not an isomorphism, there exists a monomorphism $B' \overset{j}{\hookrightarrow} B$ having the following properties:
\begin{itemize}
\item[\textup{a)}] $B' \in M_0$.
\item[\textup{b)}] $j$ does not factor through $i$.
\item[\textup{c)}] $A' = A \times_B B'$ exists in $\M$, as does $C = A \amalg_{A'} B'$, $A'$ is in $M_0$, and in the canonical diagram
\[
\begin{tikzcd}
A' \arrow[r, hook, "i'"] \arrow[d, hook] & B' \arrow[d, hook] \arrow[dr, hook, dotted, "j"] & \\
A \arrow[r, hook, "\alpha"] \arrow[rr, hook, dotted, bend right=20, "i"'] & C \arrow[r, "\beta"] & B
\end{tikzcd}
\tag*{(7.1.1)}
\]
the arrow $\beta$ is a monomorphism.
\end{itemize}
\end{quote}

(\textbf{N.B.} These conditions are satisfied if we assume, in addition to the condition that $M_0$ be generating for monomorphisms, that $M_0$ contains, along with every object $B'$, and for every object $A'$ which is either a quotient, or a subobject of $B'$, an object $A'_0$ isomorphic to $A'$; and if moreover $\M$ satisfies the following exactness conditions:
\begin{itemize}
\item[1$^{\circ}$)] Every arrow $f$ of $\M$ factors as $ip$, with $i$ a monomorphism, $p$ an epimorphism, and:
\item[2$^{\circ}$)] If $A_1$, $A_2$ are two subobjects of an object $B$ of $\M$, then the intersection $A_0 = A_1 \times_B A_2$ exists, as does the
\printedpage{47}
% head: A. Grothendieck : Dérivateurs, Ch. XV, Édition des Universités de Montpellier II et Paris VII
\msp{page 87}
amalgamated sum $C = A_1 \amalg_{A_0} A_2$, and the canonical morphism $C \longrightarrow B$ is a monomorphism.)
\end{itemize}

\textit{Assume moreover that $\M$ satisfies condition} \cercle{e}, p. 76, i.e. \textit{that every filtering $\varinjlim$ of monomorphisms is a monomorphism.}

1$^{\circ}$) \textit{Let} $\Mon \subset \Fl(\M)$ \textit{be the set of monomorphisms, and} $\Mon_0$ \textit{the subset of} $\Mon$ \textit{formed by the $i \in$} $\Mon$ \textit{such that source and target belong to $M_0$. We then have}
\[
\Mon \;\subset\; \widetilde{\Mon_0} \quad , \qquad \text{\textit{whence}} \quad
\left\{
\begin{aligned}
\widetilde{\Mon} \quad &= \quad \widetilde{\Mon_0} \\
(\Mon)_{*} \quad &= \quad (\Mon_0)_{*} \;.\text{\footnotemark[30]}
\end{aligned}
\right.
\tag*{(7.1.2)}
\]
\footnotetext[30]{\textbf{N.B.} Here we have a criterion for injectivity of objects in the categories $\M/Y$ -- it suffices to test the injectivity condition by monomorphisms $i$ which are in $\Mon_0$.}

2$^{\circ}$) \textit{If $M_0$ is} small \textit{and if the objects of $M_0$ are accessible} (e.g. $\M$ satisfies \cercle{c}, p. 76), \textit{and if} $\Mon$ \textit{is stable under cobase change, then there exists a functorial factorization of every arrow $f$ of $\M$ as}
\[
f \;=\; p\, i \;,
\]
\textit{with $i \in$} $\Mon$, \textit{$p \in$} $(\Mon)_{*}$. (\textbf{N.B.} An arrow $p : X \longrightarrow Y$ of $\M$ is in $(\Mon)_{*}$ if and only if it makes $X$ an \textit{injective} object of $\M/Y$ -- thus the factorization assertion ensures that in the induced categories $\M/Y$ ``there are enough injectives''.) \textit{Moreover, we have}
\[
\Mon \;=\; \widetilde{\Mon_0} \quad .
\tag*{(7.1.3)}
\]

3$^{\circ}$) \textit{Let $W$ be a subset of} $\Fl(\M)$ \textit{which is} L-accessible (i.e. \textit{accessible, and stable under filtering direct limits}). (\textit{We now assume that $\M$ satisfies the conditions} (a), (b), (c), (d), (e) \textit{on page 76, whence a good theory of L-accessible subsets of} $\Ob\,(\M)$ \textit{or of} $\Fl(\M)$.) \textit{We assume, more precisely, that $W$ satisfies the conditions:}

\msp{page 88}
\[
\left\{
\begin{array}{@{\quad}l@{\quad}p{0.6\textwidth}}
\text{a)} & \textit{$W$ is L-accessible} (\textit{i.e.} stable under filtering $\varinjlim$, and accessible). \\[1.5ex]
\text{b)} & \textit{Let $\alpha$, $\beta$ be composable, with $\beta\alpha$ and $\alpha$ in $W$, then $\beta \in W$.} (\textbf{N.B.} This condition will be satisfied if $W$ is a mildly saturated localizer, and it is to this situation that we are going to apply our statement.) \\[1.5ex]
\text{c)} & $\Mon \cap W$ \textit{is stable under cobase change.}
\end{array}
\right.
\tag*{(7.1.4)}
\]
\textit{and that $\M$ satisfies the additional conditions}
\[
\left\{
\begin{array}{@{\quad}l@{\quad}p{0.6\textwidth}}
\text{e$'$)} & \textit{The filtering $\varinjlim$ commute with base change.} \\[1.5ex]
\text{f)} & \textit{The exactness condition} (7.2.1) \textit{below.}
\end{array}
\right.
\tag*{(7.1.5)}
\]
\textit{Let $\pi$ be a cardinal such that $W$ is $\pi$-accessible, and $M_0 \subset$} $\Filt^{\pi}$ (\textit{hence $M_0$ essentially small}). \textit{Then we have}
\[
W \cap \Mon \;\subset\; \widetilde{(W \cap \Mon)^{\pi}}
\tag*{(7.1.6)}
\]
\printedpage{48}
% head: A. Grothendieck : Dérivateurs, Ch. XV, Édition des Universités de Montpellier II et Paris VII
(\textit{where for every subset $C$ of} $\Fl(\M)$, \textit{we set} $C^{\pi} = C \cap \Fl^{\pi}(\M)$), \textit{hence}
\[
\widetilde{W \cap \Mon} \;=\; \widetilde{(W \cap \Mon)^{\pi}} \;,
\tag*{(7.1.7)}
\]
\textit{hence $\widetilde{W \cap \Mon}$ is Q-generated by a} small \textit{subset of} $W \cap \Mon$.

\textbf{Corollary 7.2.} \textit{Let $\M$ satisfy} (a), (b), (c), (d), (e) \textit{and the additional conditions:}
\[
\left\{
\begin{array}{@{\quad}l@{\quad}p{0.6\textwidth}}
\text{e$'$)} & \textit{The filtering $\varinjlim$ commute with base change.} \\[1.5ex]
\text{f)} & \textit{If $A_1$, $A_2$ are two subobjects of an object $B$ of $\M$, and $A_0 = A_1 \cap A_2$, then $A_1 \amalg_{A_0} A_2 \hookrightarrow B$ is a monomorphism}\,\footnotemark[31].
\end{array}
\right.
\tag*{(7.2.1)}
\]
\footnotetext[31]{Another way of saying it: if $C = \sup_{\Ss(B)}(A_1, A_2)$, the cartesian square
\[
\begin{tikzcd}[ampersand replacement=\&]
A_0 \arrow[r] \arrow[d] \& A_1 \arrow[d] \\
A_2 \arrow[r] \& C
\end{tikzcd}
\]
is also cocartesian.}

\textit{Let $W$ be a subset of} $\Fl(\M)$ \textit{satisfying the following conditions.}
\[
\left\{
\begin{array}{@{\quad}l@{\quad}p{0.6\textwidth}}
\text{a)} & \textit{$W$ is a mildly saturated localizer.} \\[1.5ex]
\text{b)} & \textit{$W$ is L-accessible,} i.e. \textit{it is stable under filtering $\varinjlim$, and is accessible.} \\[1.5ex]
\text{c)} & $\Mon \cap W$ \textit{is stable under cobase change,} i.e. \textit{if we have a diagram which is} cocartesian
\begin{center}
\begin{tikzcd}[ampersand replacement=\&]
A \arrow[r, hook, "i"] \arrow[d, "\alpha"'] \arrow[dr, phantom, "\scriptstyle\text{cocart.}"] \& B \arrow[d, "\beta"] \\
A' \arrow[r, "i'"] \& B'
\end{tikzcd}
\end{center}
\textit{with $i$ a monomorphism belonging to $W$, then $i'$ is a monomorphism belonging to $W$.}\,\footnotemark[32]
\end{array}
\right.
\tag*{(7.2.3)}
\]
\footnotetext[32]{\textit{Cf.} below for a fourth condition $(d)$ on $W$, which ensures that $W$ fits into a Quillen triple.}

\msp{page 89}
\textit{Under these conditions, we have}
\[
\underbrace{\Mon}_{C} \cap W \;=\; \widetilde{W \cap \Mon} \quad , \qquad \text{i.e. } \Mon \cap W \text{ \textit{is Q-closed on the left}} \;,
\tag*{(7.2.4)}
\]